\begin{conf-abstract}
{New tool for in-situ gas measurements in low flow porous media}
{C. Marion, C. Ebi, S. Bloem, M.S. Brennwald, R. Kipfer}
{Eawag}
{Gases, including noble gases, are powerful tracers in environmental science, widely used to study groundwater flow dynamics and surface–subsurface water interactions ([1], [2], [3], [4], [5], [6]). Available technologies ([7], [8]) allow such measurement set-up only for high water flow environmental systems (L/min to L/s); systems in which the extraction of liters of water per minute would not disturb the systems equilibrium. In the meanwhile, in-situ gas measurements in low-flow porous media (L/h), such as soil unsaturated and saturated zones or plant tissues, are technically challenging, even though these systems play a critical role for water and CO2 exchange on the global scale. 
We present a robust, field-deployable gas probe for long-term, in-situ gas monitoring in low-permeability or low-flow environments. Originally developed for in-vivo gas monitoring in tree xylem ([9]), this method enables non-destructive, high-resolution measurements applicable to a variety of porous media, including soil, sediment, rock, and biological tissue.

[1] Brennwald, M. S. et al. (2022). Frontiers in Water, 4, 107.
[2] Blanc, T. et al. (2024). Water Research, 254, 121375. 
[3] Schilling, O. S. et al. (2017). Water Resources Research, 53(12), 10465–10490.
[4] Schilling, O. S. et al. (2021). Water Resources Research, 57(2).
[5] Peel, M. et al. (2022). Chemical Geology, 599, 120829.
[6] Giroud, S. et al. (2023). Frontiers in Water, 4, 1032094.
[7] Brennwald, M. S., et al. (2016). Environmental Science and Technology, 50.
[8] Maechler, L. et al. (2014). Chimia, 68(3), 155–159. 
[9] Marion, C. et al. (2024). Tree Physiology. }
\end{conf-abstract}
